\documentclass[10pt,a4paper,sans]{moderncv}

\moderncvstyle{banking}
\moderncvcolor{blue}

\usepackage[utf8]{inputenc}
\usepackage[scale=0.84]{geometry}
\usepackage{xcolor}
\usepackage{hyperref}

\definecolor{cvblue}{HTML}{0A66C2}
\colorlet{color1}{cvblue}
\colorlet{color2}{black}
\colorlet{firstnamecolor}{cvblue}
\colorlet{lastnamecolor}{cvblue}
\colorlet{namecolor}{cvblue}
\colorlet{sectioncolor}{cvblue}
\colorlet{subsectioncolor}{cvblue}
\colorlet{bodyrulecolor}{cvblue}
\colorlet{headrulecolor}{cvblue}
\AtBeginDocument{\hypersetup{colorlinks=true,urlcolor=cvblue,linkcolor=cvblue}}


\name{Thiago}{Costa}
\address{Teresópolis, RJ, BR}{}{}
%\phone[mobile]{+1 561 808 6390}
\phone[mobile]{+55 24 98836 7663}
\email{tabrantedaco2023@fau.edu}
\social[linkedin]{thiagosieg}
\social[instagram]{sieg.ia}

\begin{document}
\makecvtitle

\section{RESUMO PROFISSIONAL}
\cvitem{}{Mestrando em Inteligência Artificial na Florida Atlantic University, com experiência em aplicações de IA, machine learning, sistemas multiagentes, RAG, prototipação de produtos e experimentação orientada por dados. Atua na construção de soluções com Python, SQL, PyTorch, scikit learn, LangChain e bancos vetoriais, conectando pesquisa aplicada, engenharia e visão de produto. Possui perfil analítico, comunicação técnica forte e experiência em transformar problemas complexos em soluções práticas. Interesses de pesquisa: orquestração multiagente com LLMs e quantificação de incerteza em modelos de deep learning para séries temporais. Reconhecido com segundo lugar no Agentic AI Hackathon da FAU em 2025 e no Three Minutes International Speech Competition em 2024.}

\section{FORMAÇÃO ACADÊMICA}
\cventry{2024 a 2026}{Mestrado em Inteligência Artificial}{Florida Atlantic University}{Boca Raton, FL, EUA}{}{}
\cventry{2020 a 2022}{Bacharelado em Engenharia Mecânica}{Universidade Federal Fluminense}{Volta Redonda, RJ, Brasil}{}{}
\cventry{2017 a 2020}{Bacharelado em Física Computacional}{Universidade Federal Fluminense}{Volta Redonda, RJ, Brasil}{Curso não concluído}{}

\section{COMPETÊNCIAS TÉCNICAS}
\cvitem{Linguagens}{Python, SQL, C, C++, HTML, LaTeX}
\cvitem{IA e ML}{Machine Learning supervisionado e não supervisionado, Deep Learning, LLMs, RAG, sistemas multiagentes, NLP, avaliação de modelos, engenharia de prompts}
\cvitem{Bibliotecas}{PyTorch, scikit learn, pandas, numpy, matplotlib, LangChain}
\cvitem{Dados}{Análise exploratória, visualização de dados, preparação de dados, métricas, experimentação e validação}
\cvitem{Ferramentas}{Git, GitHub, bancos vetoriais, Ollama, Chroma, FAISS, desenvolvimento de protótipos}
\cvitem{Idiomas}{Português nativo, Inglês fluente, Espanhol intermediário}

\section{EXPERIÊNCIA RELEVANTE}
\cventry{Novembro 2025 a atual}{AI Engineer}{NET, Neurodiverse Educational Testing}{}{Startup}{ }
\cvitem{}{Desenvolveu a Web UI e integrou uma arquitetura multiagente para triagem educacional assistida por IA.}
\cvitem{}{Implementou um agente conversacional com guardrails para coleta estruturada de contexto, um pipeline RAG com retrieval semântico e matching com base de conhecimento sobre transtornos de aprendizagem, e um agente gerador de relatórios em PDF com recomendações e encaminhamentos para avaliação especializada.}
\cvitem{}{Conectou interface, orquestração e fluxo de saída do sistema em um protótipo funcional. Projeto reconhecido com segundo lugar no FAU Agentic AI Hackathon de 2025.}

\vspace{5 mm}

\cventry{Maio 2025 a atual}{Criador e Gestor}{SIEG IA}{}{Canal de conteúdo sobre Inteligência Artificial}{ }
\cvitem{}{Construiu e escalou um canal de conteúdo focado em notícias e tendências de IA, alcançando mais de 100 mil seguidores e 100 milhões de visualizações em 2025.}
\cvitem{}{Transforma temas técnicos complexos em comunicação acessível, fortalecendo autoridade em IA, leitura de tendências e posicionamento de mercado.}



\section{PROJETOS SELECIONADOS}
\cventry{\textbf{Projeto de IA aplicada}}{\textbf{Sistema multiagente com RAG para apoio neuroeducacional}}{}{}{}{}
\cvitem{}{Estruturou um fluxo com agente conversacional protegido por guardrails, recuperação semântica com RAG e geração automatizada de relatórios em PDF para apoio a encaminhamentos educacionais e clínicos.}
\cvitem{}{Trabalhou na integração entre interface Web, orquestração entre agentes e entrega final orientada ao usuário.}

\cventry{\textbf{Projeto de Deep Learning}}{\textbf{Modelo convolucional para xadrez}}{}{}{}{}
\cvitem{}{Desenvolveu modelo em PyTorch para processar representações do tabuleiro em 13 canais, utilizando camadas convolucionais e camadas densas para aprender padrões táticos e posicionais em xadrez.}
\cvitem{}{Aplicou arquitetura de deep learning para classificação e suporte à tomada de decisão em cenários do jogo.}

\cventry{\textbf{Projeto de Machine Learning}}{\textbf{Previsão de preços com LSTM}}{}{}{}{}
\cvitem{}{Desenvolveu modelo para previsão de preços de criptoativos com séries temporais, avaliando modelagem sequencial e potencial uso de indicadores técnicos.}

\cventry{\textbf{Projeto de Reinforcement Learning}}{\textbf{Agente para ambiente tabular e grid world}}{}{}{}{}
\cvitem{}{Implementou soluções com aprendizado por reforço e Deep Q Network em ambientes de decisão sequencial, utilizando PyTorch para treinamento e avaliação.}

\section{DIFERENCIAIS}
\cvitem{}{Mentalidade de produto, forte autonomia de execução, facilidade para aprender novas ferramentas rapidamente e boa comunicação entre áreas técnicas e não técnicas.}
\cvitem{}{Experiência internacional, atuação em ambiente acadêmico nos Estados Unidos e produção contínua em temas avançados de Inteligência Artificial.}

\section{LIDERANÇA E ATIVIDADES}
\cvitem{}{Líder de Projeto, Equipe de Lançamento de Foguetes da UFF.}
\cvitem{}{Monitor de Mecânica dos Fluidos na UFF.}
\cvitem{}{Experiência complementar em liderança estudantil, representação acadêmica e comunicação técnica.}

\section{PRÊMIOS E RECONHECIMENTOS}
\cventry{2025}{Segundo lugar, Agentic AI Hackathon}{Florida Atlantic University}{}{}{Reconhecimento pelo projeto NET, solução com agentes de IA voltada a apoiar estudantes neurodivergentes e suas famílias.}
\cventry{2024}{Segundo lugar, Three Minutes International Speech Competition}{}{}{}{Reconhecimento por apresentação em formato pitch sobre o impacto da Inteligência Artificial no dia a dia.}

\end{document}
